\input{../tex-inputs/latex-leseansicht-vorspann}

\section[Arthur und Olga Schnitzler an Gustav Schwarzkopf, 15. 5. 1914]{L04486 Arthur und Olga Schnitzler an Gustav Schwarzkopf, 15. 5. 1914}
\nopagebreak\mylabel{L04486v}
\rehead{ }\normalsize\beginnumbering\briefempfaengerindex{Schwarzkopf, Gustav@\textsc{Schwarzkopf, Gustav}!zzzSchnitzler, Olga@\emph{von Olga Schnitzler}!1914-05-151@{15. 5. 1914}|(be}\briefempfaengerindex{Schwarzkopf, Gustav@\textsc{Schwarzkopf, Gustav}!zzzSchnitzler, Arthur@\emph{von Arthur Schnitzler}!1914-05-151@{15. 5. 1914}|(be}
\toendnotes[C]{\smallbreak\pagebreak[2]}
\correspDesc{Versand  durch Arthur Schnitzler, Olga Schnitzler am 15. 5. 1914 in Algier
\newline{}Übermittlung  am 15. 5. 1914 in Algier
\newline{}Erhalt  durch Gustav Schwarzkopf im Zeitraum [16. 5. 1914 – 20. 5. 1914?] in Wien}\toendnotes[C]{\smallbreak}
\Standort{DLA, A:Schnitzler, HS.1985.1.1897.}
\physDesc{Bildpostkarte, 243 Zeichen
\newline{}Handschrift Arthur Schnitzler: Bleistift, deutsche Kurrent
\newline{}Handschrift Olga Schnitzler: Bleistift, lateinische Kurrent
\newline{}Versand: Stempel: »\nobreak{}\oindex{Algier@\textbf{Algier}|pwk}Alger / R. de Strasbourg, 15–5, 1\textcolor{gray}{6}{[}.{]}45\nobreak{}«.  }\toendnotes[C]{\smallbreak}\pstart{}{\pb}\textsc{Austria\oindex{Österreich@\textbf{Österreich}|pw}}\pend{}\pstart{}Herrn \textsc{Gustav Schwarzkopf}\pend{}\pstart{}Wien I\oindex{I., Innere Stadt@\textbf{I., Innere Stadt}|pw}\pend{}\pstart{}\textsc{Tiefer Graben 17}\oindex{Wien@\textbf{Wien}!I., Innere Stadt@\textbf{I., Innere Stadt}!Tiefer Graben 17@\textbf{Tiefer Graben 17}|pw}\pend{}{\bigskip}
\pstart
           \noindent{}\centering{}{\pb}\textcolor{gray}{\textbf{1314. ALGER. – Quais – Square Bresson et Tournants Rovigo\oindex{Place Port Saïd@\textbf{Place Port Saïd}|pw} – E. S.}}\pend
           \vspace{1em}
\pstart
           {\pb}\textsc{Algier\oindex{Algier@\textbf{Algier}|pw}}{ }15/5/914 –\pend
           \vspace{0.5em}
\pstart
           Herzliche Grüße.\pend
           
\pstart
           Wir{ }ſind 5 Stunden hier gewesen – fahren 
      nun weiter nach
            {\pb}Liſſabon\oindex{Lissabon@\textbf{Lissabon}|pw}, die Fahrt iſt bisher ſo{ }ſchön als wie möglich! Das Schiff\orgindex{Reichspostdampfer Yorck@Reichspostdampfer Yorck|pwv}
      höchſt angenehm{[}.{]} Ihr \spacefill\mbox{A.}\pend
           \selectlanguage{ngerman}\vspace{1em}
\pstart
           \noindent{}{[}hs. O. Schnitzler:{]} Herzlichst \spacefill\mbox{Olga.}\pend
           \selectlanguage{ngerman}\endnumbering\briefempfaengerindex{Schwarzkopf, Gustav@\textsc{Schwarzkopf, Gustav}!zzzSchnitzler, Olga@\emph{von Olga Schnitzler}!1914-05-151@{15. 5. 1914}|)be}\briefempfaengerindex{Schwarzkopf, Gustav@\textsc{Schwarzkopf, Gustav}!zzzSchnitzler, Arthur@\emph{von Arthur Schnitzler}!1914-05-151@{15. 5. 1914}|)be}\mylabel{L04486h}  \newcommand{\dateiname}{L04486}\newcommand{\titel}{Arthur und Olga Schnitzler an Gustav Schwarzkopf, 15. 5. 1914}\newcommand{\editorInnen}{Herausgegeben von Jahnke, SelmaMüller, Martin Anton}\input{../tex-inputs/latex-leseansicht-abspann}
